\section{Autentifikavimas MFA metodu}
\label{sec:autentifikavimas}

\subsection{Prisijungimo eiga}
\label{subsec:prisijungimas}

Vartotojas, kuriam MFA aktyvuotas, prisijungia dviem žingsniais (\ref{fig:prisijungimas} pav.):
\begin{enumerate}
    \item Vartotojas įveda el. pašto adresą ir slaptažodį (pirmasis faktorius).
    \item Aplikacija patikrina slaptažodžio maišą (argon2id). Jei slaptažodis teisingas ir \texttt{mfa\_enabled = true}, sesija dar \emph{nesukuriama}. Išsaugoma tik laukiančioji būsena (\textit{angl. pending}): \texttt{pending\_user\_id}, \texttt{state}, \texttt{nonce} ir \texttt{code\_verifier}. Ji galioja 5 min.
    \item Naršyklė peradresuojama į \textit{Google}. Parametro \texttt{login\_hint} reikšmė yra susietos paskyros \texttt{sub}, todėl \textit{Google} iš karto pasiūlo teisingą paskyrą~\cite{google_oidc}.
    \item Vartotojas prisijungia prie \textit{Google}. Jei naršyklėje jau yra aktyvi \textit{Google} sesija, pakanka ją patvirtinti.
    \item \textit{Google} grąžina vartotoją į \texttt{redirect\_uri}. Aplikacija patikrina \texttt{state}, iškeičia kodą į žetonus (su \texttt{code\_verifier}), patikrina ID žetoną ir \texttt{nonce}. Tada \textbf{palyginamas ID žetono \texttt{sub} su susietu \texttt{provider\_sub}}. Produkcinėje aplikacijoje papildomai tikrinama, ar \texttt{amr} turi \texttt{mfa}, \texttt{hwk} arba \texttt{swk} ir ar \texttt{auth\_time} pakankamai naujas.
    \item \textbf{Sėkmė.} Sukuriama nauja sesija su nauju sesijos identifikatoriumi, kad būtų išvengta sesijos fiksavimo (\textit{angl. session fixation}). Atnaujinamas \texttt{last\_used\_at}, laukiančioji būsena ištrinama.
    \item \textbf{Nesėkmė} (vartotojas atšaukė prisijungimą, \texttt{sub} nesutampa, žetonas netinkamas arba baigėsi laukimo laikas). Prisijungimas atmetamas, įvykis įrašomas į \texttt{auth\_events}, bandymų skaičius ribojamas. Vartotojui išsiunčiamas pranešimas su bandymo laiku, naršykle ir vieta~\cite{VMA}, o prisijungimo lange pasiūloma naudoti atkūrimo kodą.
\end{enumerate}

\begin{figure}[H]
    \centering
    \begin{adjustbox}{max width=\textwidth}
    \begin{tikzpicture}
        % Dalyviai
        \node[actor] (U) at (0,0)  {Vartotojas\\(naršyklė)};
        \node[actor] (A) at (5,0)  {Aplikacija\\(serveris)};
        \node[actor] (G) at (10,0) {\textit{Google}\\(OpenID teikėjas)};
        \node[actor] (D) at (15,0) {Duomenų\\bazė};

        % Gyvavimo linijos
        \foreach \x in {0,5,10,15} {
            \draw[lifeline] (\x,-0.5) -- (\x,-23.0);
        }

        \draw[msg] (0,-1.4) -- node[lbl, above] {1. El. paštas + slaptažodis} (5,-1.4);

        \draw[msg] (5,-2.4) -- node[lbl, above, pos=0.75] {2. Slaptažodžio patikra (argon2id)} (15,-2.4);

        \draw[ret] (15,-3.4) -- node[lbl, above, pos=0.25] {2. \texttt{mfa\_enabled}, \texttt{provider\_sub}} (5,-3.4);

        \node[note] at (5,-4.7) {2. Slaptažodis teisingas $\to$ tik laukiančioji būsena:\\
            \texttt{pending\_user\_id}, \texttt{state}, \texttt{nonce},\\
            \texttt{code\_verifier} (galioja 5 min.)};

        \draw[ret] (5,-6.3) -- node[lbl, above] {3. HTTP 302 $\to$ \textit{Google}} (0,-6.3);

        \draw[msg] (0,-7.5) -- node[lbl, above, pos=0.25] {3. \texttt{login\_hint} = susietas \texttt{sub}} (10,-7.5);

        \node[note] at (10,-8.7) {4. Prisijungimas prie \textit{Google}\\
            (jei sesija aktyvi -- tik patvirtinimas)};

        \draw[ret] (10,-10.4) -- node[lbl, above, pos=0.75] {5. HTTP 302 $\to$ \texttt{redirect\_uri}\\
            (\texttt{code}, \texttt{state})} (0,-10.4);

        \draw[msg] (0,-11.6) -- node[lbl, above] {5. GET \texttt{/auth/google/callback}} (5,-11.6);

        \draw[msg] (5,-12.8) -- node[lbl, above] {5. POST \texttt{/token} (+\texttt{code\_verifier})} (10,-12.8);

        \draw[ret] (10,-13.8) -- node[lbl, above] {5. \texttt{id\_token}} (5,-13.8);

        \node[note] at (5,-15.1) {5. Tikrina \texttt{state}, ID žetoną, \texttt{nonce};\\
            \textbf{ar \texttt{sub} = susietas \texttt{provider\_sub}?}};

        % alt fragmentas: sėkmė arba nesėkmė
        \draw[thick] (-1.9,-16.2) rectangle (16.9,-22.6);
        \node[draw, thick, fill=white, font=\footnotesize\bfseries, anchor=north west,
            inner sep=2pt] at (-1.9,-16.2) {alt};
        \node[font=\footnotesize, fill=white, anchor=north west, inner sep=1.5pt]
            at (-1.0,-16.25) {[\texttt{sub} sutampa]};

        \draw[msg] (5,-17.3) -- node[lbl, above, pos=0.75] {6. \texttt{last\_used\_at}, \texttt{mfa\_success}} (15,-17.3);

        \draw[ret] (5,-18.3) -- node[lbl, above] {6. Nauja sesija (naujas ID)} (0,-18.3);

        \draw[dashed] (-1.9,-18.9) -- (16.9,-18.9);
        \node[font=\footnotesize, fill=white, anchor=north west, inner sep=1.5pt]
            at (-1.8,-18.95) {[\texttt{sub} nesutampa / atšaukta / klaida]};

        \draw[msg] (5,-20.0) -- node[lbl, above, pos=0.75] {7. \texttt{auth\_events}: \texttt{mfa\_failed}} (15,-20.0);

        \draw[ret] (5,-21.0) -- node[lbl, above] {7. Prisijungimas atmestas} (0,-21.0);

        \draw[msg] (5,-22.0) -- node[lbl, above] {7. Įspėjimas el. paštu} (0,-22.0);
    \end{tikzpicture}
    \end{adjustbox}
    \caption{Prisijungimas su MFA: slaptažodis ir \textit{Google} patvirtinimas.}
    \label{fig:prisijungimas}
\end{figure}


\subsection{Kada dar reikalauti MFA}
\label{subsec:step_up}

Pagal skaidrių rekomendacijas~\cite{VMA} OAuth patvirtinimas reikalingas ne tik prisijungiant. Tas pats srautas (su \texttt{purpose = step\_up}) naudojamas ir kai vartotojas:
\begin{itemize}
    \item keičia slaptažodį arba el. pašto adresą;
    \item išjungia MFA arba keičia susietą \textit{Google} paskyrą;
    \item sesijoje gauna aukštesnes teises (pvz., administratoriaus).
\end{itemize}
Tokiems veiksmams priimamas tik neseniai atliktas patvirtinimas (pvz., per paskutines 5 min.). Be to, MFA turi būti taikomas visiems prisijungimo būdams. Pavyzdžiui, slaptažodžio atkūrimas el. paštu neturi leisti apeiti \textit{Google} žingsnio.

\subsection{Vartotojo patirties gerinimas}
\label{subsec:rizika}

Dažnas MFA reikalavimas gali erzinti vartotojus, todėl galima taikyti rizika paremtą MFA~\cite{VMA}. Sėkmingai prisijungus, įrenginiui išduodamas pasirašytas „patikimo įrenginio“ slapukas (\texttt{HttpOnly}, \texttt{Secure}, galioja 30 d.). Tada \textit{Google} žingsnio reikalaujama tik jungiantis iš naujo įrenginio ar neįprastos vietos arba po nesėkmingų bandymų.

\subsection{MFA atkūrimas}
\label{subsec:atkurimas}

Praradęs prieigą prie \textit{Google} paskyros, vartotojas negalėtų prisijungti. Todėl numatomi atkūrimo būdai~\cite{VMA}:
\begin{itemize}
    \item vienkartiniai atkūrimo kodai, gauti aktyvuojant MFA. Panaudotas kodas pažymimas, o vartotojas raginamas iš naujo susieti \textit{Google} paskyrą;
    \item galimybė susieti antrą MFA metodą (pvz., TOTP programėlę), kad vieno metodo praradimas neužrakintų paskyros;
    \item kraštutiniu atveju -- kreipimasis į klientų aptarnavimą, kuris griežtai patikrina vartotojo tapatybę.
\end{itemize}
